%%
%% This is file `main.tex' based on `sample-sigconf.tex' (q.v. for spurce of that,
%%
%% IMPORTANT NOTICE:
%% 
%% For the copyright see the original source file `sample-sigconf.tex'
%% in the `Sample' folder.
%%
%% For distribution of the original source see the terms
%% for copying and modification in the file samples.dtx.
%% 
%% This generated file may be distributed as long as the
%% original source files, as listed above, are part of the
%% same distribution. (The sources need not necessarily be
%% in the same archive or directory.)
%%
%% Commands for TeXCount
%TC:macro \cite [option:text,text]
%TC:macro \citep [option:text,text]
%TC:macro \citet [option:text,text]
%TC:envir table 0 1
%TC:envir table* 0 1
%TC:envir tabular [ignore] word
%TC:envir displaymath 0 word
%TC:envir math 0 word
%TC:envir comment 0 0
%%
%%
%% The first command in your LaTeX source must be the \documentclass command.

%% NOTE that a single column version is required for 
%% submission and peer review. This can be done by changing
%% the \doucmentclass[...]{acmart} in this template to 
%%\documentclass[manuscript,review,anonymous]{acmart}
%% This version is used for drafting and final submission
\documentclass[sigconf, 10pt]{acmart}
\settopmatter{printacmref=false}

\usepackage{url}


%% 
%% To ensure 100% compatibility, please check the white list of
%% approved LaTeX packages to be used with the Master Article Template at
%% https://www.acm.org/publications/taps/whitelist-of-latex-packages 
%% before creating your document. The white list page provides 
%% information on how to submit additional LaTeX packages for 
%% review and adoption.
%% Fonts used in the template cannot be substituted; margin 
%% adjustments are not allowed.

%%
%% \BibTeX command to typeset BibTeX logo in the docs
\AtBeginDocument{%
  \providecommand\BibTeX{{%
    \normalfont B\kern-0.5em{\scshape i\kern-0.25em b}\kern-0.8em\TeX}}}

%% The following are not a requirement, delete if not using
\acmDOI{}
\acmISBN{}
\setcopyright{none}
\acmConference[ACL]{}{Algorithms for Computational Logic}{19-10-2025}

\makeatletter
\patchcmd{\maketitle}{\@copyrightyear.}{}{}{}
\makeatother

%%
%% Submission ID.
%% Use this when submitting an article to a sponsored event. You'll
%% receive a unique submission ID from the organizers
%% of the event, and this ID should be used as the parameter to this command.
%%\acmSubmissionID{123-A56-BU3}

%%
%% For managing citations, it is recommended to use bibliography
%% files in BibTeX format.
%%
%% You can then either use BibTeX with the ACM-Reference-Format style,
%% or BibLaTeX with the acmnumeric or acmauthoryear sytles, that include
%% support for advanced citation of software artefact from the
%% biblatex-software package, also separately available on CTAN.
%%
%% Look at the sample-*-biblatex.tex files for templates showcasing
%% the biblatex styles.
%%

%%
%% The majority of ACM publications use numbered citations and
%% references.  The command \citestyle{authoryear} switches to the
%% "author year" style.
%%
%% If you are preparing content for an event
%% sponsored by ACM SIGGRAPH, you must use the "author year" style of
%% citations and references.
%% Uncommenting
%% the next command will enable that style.
%%\citestyle{acmauthoryear}

%%
%% end of the preamble, start of the body of the document source.

%% % Location of your graphics files for figures, here a sub-folder to the main project folder
\graphicspath{{./images/}} 


\begin{document}

%%
%% The "title" command has an optional parameter,
%% allowing the author to define a "short title" to be used in page headers.
\title{The Nurse Scheduling Problem}

%%
%% The "author" command and its associated commands are used to define
%% the authors and their affiliations.
%% Of note is the shared affiliation of the first two authors, and the
%% "authornote" and "authornotemark" commands
%% used to denote shared contribution to the research.
\author{Francisco Fonseca, 102492}
\authornote{Both authors contributed equally to this work.}
\email{franciscodiasfonseca@tecnico.ulisboa.pt}
\author{Gonçalo Rua, 102604}
\authornotemark[1]
\email{goncalonrua@tecnico.ulisboa.pt}
\affiliation{%
  \institution{Instituto Superior Técnico}
  \streetaddress{Av. Rovisco Pais}
  \city{Lisbon}
  \country{Portugal}
}

%%
%% By default, the full list of authors will be used in the page
%% headers. Often, this list is too long, and will overlap
%% other information printed in the page headers. This command allows
%% the author to define a more concise list
%% of authors' names for this purpose.
\renewcommand{\shortauthors}{Francisco Fonseca, Gonçalo Rua}

%%
%% This command processes the author and affiliation and title
%% information and builds the first part of the formatted document.
\maketitle

\section{Description of the Problem}
The Nurse Scheduling Problem (NSP) involves creating an optimal shift schedule for nurses working in a hospital over a given period of days, subject to some constraints:

\begin{enumerate}
    \item Each shift requires a minimum number of nurses for the normal functioning of the hospital.
    \item Each shift must have at least one nurse with managerial qualifications.
    \item A nurse's assigned shifts on any given day must be continuous, with no gaps in between.
    \item Each nurse has a minimum and maximum daily workload, defined by the number of assigned shifts.
    \item The total number of shifts that can be assigned to a nurse for the entire period is limited.
    \item A nurse cannot be assigned to shifts when it's known they're unavailable to work.
    \item The schedule must consider nurses' shift preferences. An optimal solution aims to minimize the total cost associated with the assignments.
\end{enumerate}

\section{Description of the Encoding}

We first define a set of variables to represent the schedule and then encode the problem's constraints as clauses.

\subsection{Definition of Variables}
\begin{itemize}
    \item Let $D$ be a positive integer representing the number of days in the scheduling period.
    \item Let $S$ be a positive integer representing the number of shifts per day, all with the same duration, that cover the full 24 hours. For example, if $S = 4$, each shift will have 6 hours.
    \item Let $R_{s,d}$ be the minimum number of nurses required for shift $s$ ($1 \le s \le S$) on day $d$ ($1 \le d \le D$).
    \item Let $N = \{N_1, N_2, \dots, N_n\}$ be the set of $n$ available nurses.
    \item Let $M \subseteq N$ be the subset of nurses with managerial qualifications.
    \item Let $k_{\min}$ and $k_{\max}$ be the minimum and maximum number of shifts for nurse $N_k$ ($1 \le k \le n$) in any given day.
    \item Let $H_k$ be the maximum total number of shifts that can be assigned to nurse $N_k$ over the entire period.
    \item Let $P_{k,s,d}$ be the cost of assigning nurse $N_k$ to shift $s$ on day $d$. If $P_{k,s,d} = -1$, nurse $N_k$ is not available for that shift. All other costs are positive integers representing preferences (a lower cost is better).
\end{itemize}

\subsection{Modeling Constraints}

This problem can be modeled and solved using SMT. 

We define a set of integer variables that describe the schedule more abstractly than using boolean assignments. For each nurse $k \in \{1, \dots, n\}$ and each day $d \in \{0, \dots, D-1\}$, we define two core integer variables:
\begin{itemize}
    \item $\mathbf{start_{k,d}}$: The 0-indexed starting shift for nurse $k$ on day $d$.
    \item $\mathbf{num_{k,d}}$: The total number of consecutive shifts nurse $k$ works on day $d$.
\end{itemize}

To simplify the main constraints, we first define a helper predicate, $Works(k, d, s)$, which is true if nurse $k$ is working on shift $s$ of day $d$. The shift index $s$ is 0-indexed, $s \in \{0, \dots, S-1\}$.
\begin{equation}
    Works(k, d, s) \equiv (start_{k,d} \le s) \land (s < start_{k,d} + num_{k,d})
\end{equation}

The constraints are then defined as follows, using an indicator function $\mathbb{I}(\cdot)$ which is 1 if its argument is true and 0 otherwise. This corresponds to the \textit{If(condition, 1, 0)} construct in the code.

\begin{enumerate}
    \item \textbf{Domain Constraints:} The variables must stay within physically possible bounds. For all $k \in \{1, \dots, n\}$ and $d \in \{0, \dots, D-1\}$:
    \begin{gather}
        0 \le start_{k,d} < S \\
        0 \le num_{k,d} \le S \\
        start_{k,d} + num_{k,d} \le S
    \end{gather}

    \item \textbf{Minimum Nurses per Shift:}
    \begin{equation}
    \begin{gathered}
        \sum_{k=1}^{n} \mathbb{I}(Works(k, d, s)) \ge R_{s,d}, \\
        \forall d \in \{0, \dots, D-1\}, s \in \{0, \dots, S-1\}
    \end{gathered}
    \end{equation}

    \item \textbf{Manager Requirement:}
    \begin{equation}
    \begin{gathered}
        \sum_{k \in M} \mathbb{I}(Works(k, d, s)) \ge 1, \\
        \forall d \in \{0, \dots, D-1\}, s \in \{0, \dots, S-1\}
    \end{gathered}
    \end{equation}

    \item \textbf{Shift Continuity:} This constraint is satisfied implicitly by the variable encoding. By defining a nurse's daily work as a single block of $num_{k,d}$ shifts starting at $start_{k,d}$, gaps within a work period are impossible. Therefore, no separate constraint is needed to enforce this rule.

    \item \textbf{Daily Workload:}
    \begin{equation}
    \begin{gathered}
        (L_{k}^{min} \le num_{k,d} \le L_{k}^{max}), \\
        \forall k \in \{1, \dots, n\}, d \in \{0, \dots, D-1\}
    \end{gathered}
    \end{equation}

    \item \textbf{Total Workload:}
    \begin{equation}
        \sum_{d=0}^{D-1} num_{k,d} \le H_k, \quad \forall k \in \{1, \dots, n\}
    \end{equation}

    \item \textbf{Unavailability:}
    \begin{equation}
        \neg Works(k, d, s), \quad \forall k, d, s \text { where } P_{k,s,d} = -1
    \end{equation}

    \item \textbf{Objective Function:}
    \begin{equation}
        \text{minimize } \sum_{d=0}^{D-1} \sum_{s=0}^{S-1} \sum_{k=1, P_{k,s,d} \ge 0}^{n} \mathbb{I}(Works(k, d, s)) \cdot P_{k,s,d}
    \end{equation}
\end{enumerate}

\section{Alternative Encodings}

\subsection{Boolean Encoding}

Let $x_{k,s,d}$ be a Boolean variable such that $x_{k,s,d} = 1$ if nurse $N_k$ is assigned to shift $s$ on day $d$, and $x_{k,s,d} = 0$ otherwise.

\begin{enumerate}
    \item \textbf{Minimum nurses per shift:}
    \begin{equation}
        \sum_{k=1}^{n} x_{k,s,d} \geq R_{s,d}, \quad \forall s \in [1,S], d \in [1,D]
    \end{equation}

    \item \textbf{Manager requirement:}
    \begin{equation}
        \sum_{N_k \in M} x_{k,s,d} \geq 1, \quad \forall s \in [1,S], d \in [1,D]
    \end{equation}

    \item \textbf{Shift continuity:}
    \begin{equation}
    \begin{gathered}
        ( \neg x_{k,s,d} \lor x_{k,s+1,d} \lor \neg x_{k,v,d} ), \\
        \forall N_k \in N, d \in [1,D], s \in [1,S-1], v \in [s+2,S]
    \end{gathered}
    \end{equation}

    \item \textbf{Daily workload:}
    \begin{equation}
        k_{\min} \leq \sum_{s=1}^{S} x_{k,s,d} \leq k_{\max}, \quad \forall N_k \in N, d \in [1,D]
    \end{equation}

    \item \textbf{Total workload:}
    \begin{equation}
        \sum_{d=1}^{D} \sum_{s=1}^{S} x_{k,s,d} \leq H_k, \quad \forall N_k \in N
    \end{equation}

    \item \textbf{Unavailability:}
    \begin{equation}
        (\neg x_{k,s,d}), \quad \forall k,s,d \text{ where } P_{k,s,d} = -1
    \end{equation}

    \item \textbf{Objective Function:}
    \begin{equation}
        \text{minimize } \sum_{d=0}^{D-1} \sum_{s=0}^{S-1} \sum_{k=1, P_{k,s,d} \ge 0}^{n} x_{k,d,s} \cdot P_{k,s,d}
    \end{equation}
\end{enumerate}

\subsection{Hybrid Encoding}

The hybrid encoding seeks to combine the strengths of both the pure boolean and pure integer models. It uses integer variables to elegantly handle structural constraints like shift continuity and workload totals, while also using boolean variables to allow for fast, simple checks on a per-shift basis.

This model defines three sets of variables for each nurse $k$, day $d$, and shift $s$:
\begin{itemize}
    \item A boolean variable $x_{k,d,s}$, which is true if nurse $k$ is assigned to shift $s$ on day $d$.
    \item An integer variable $start_{k,d}$ for the starting shift of nurse $k$ on day $d$.
    \item An integer variable $num_{k,d}$ for the number of consecutive shifts worked by nurse $k$ on day $d$.
\end{itemize}

The core of this model is a "bridge" constraint that links the boolean and integer variables, ensuring they are consistent. For all $k, d, s$:
\begin{equation}
    x_{k,d,s} \iff (start_{k,d} \le s) \land (s < start_{k,d} + num_{k,d})
\end{equation}

With this bridge in place, we can formulate each constraint using the most efficient variable type for the task:

\begin{enumerate}
    \item \textbf{Per-Shift Constraints:} The requirements for \textbf{minimum nurses per shift}, \textbf{managers requirement}, and \textbf{unavailability} are best expressed using the simple boolean variables. These constraints are identical to the pure boolean encoding.
    
    \item \textbf{Structural Constraints:} The rules for \textbf{daily workload} and \textbf{total workload} are handled by the more expressive integer variables. These constraints are identical to the first encoding. For \textbf{shift continuity}, it is satisfied by the bridge constraint linking $x_{k,d,s}$ to the single, continuous block defined by $start_{k,d}$ and $num_{k,d}$.
    
    \item \textbf{Objective Function:} The cost function is also defined using the simple boolean variables, summing the preferences for all assigned shifts.
    \begin{equation}
        \text{minimize } \sum_{d=0}^{D-1} \sum_{s=0}^{S-1} \sum_{k=1, P_{k,s,d} \ge 0}^{n} x_{k,d,s} \cdot P_{k,s,d}
    \end{equation}
\end{enumerate}

\section{Performance Comparison}

TODO

\section{Algorithm and Configuration}
The encoding was implemented in Python, and the resulting formula was provided to Microsoft's Z3 SMT Solver \cite{z3}, accessed via the \texttt{z3-solver} library. The \texttt{Optimize} class is used to handle the problem, as it is designed to find a model that not only satisfies all the hard constraints but also minimizes an objective function, which in this case is the total assignment cost. We set the solver strategy to use a Pareto approach (via \texttt{solver.set(priority='pareto')}) since we empirically observe that it provides a significant improvement in solving time compared to the default.

\section{Installing and running}
This project was implemented in Python. The sole dependency is z3-python, which can be installed via pip:
\begin{verbatim}
pip3 install z3-python
\end{verbatim}
The program can be executed with the following command:
\begin{verbatim}
python3 project2.py < [input.txt] > [output.txt]
\end{verbatim}

%%
%% The next two lines define the bibliography style to be used, and
%% the bibliography file.
\bibliographystyle{ACM-Reference-Format}
\bibliography{references.bib}

\newpage

\end{document}
